\chapter{Δεδομένα Ενεργειακής Κατανάλωσης και Αξιολόγησης}

\noindent Τα δεδομένα του παρόντος παραρτήματος παρουσιάζονται ως περιγραφική στατιστική της περιόδου παρακολούθησης. Η παρατηρηθείσα μείωση κατανάλωσης συσχετίζεται με την εγκατάσταση του συστήματος αλλά δεν αποδεικνύει αιτιακή σχέση λόγω απουσίας υποδιαίρεσης κατανάλωσης ανά υποσύστημα.

Τα παρακάτω δεδομένα προέρχονται από τους επίσημους λογαριασμούς ηλεκτρικής ενέργειας (\texttt{Τριφασική παροχή 135 \textlatin{kVA}}) της εγκατάστασης. Η σύγκριση πραγματοποιείται μεταξύ της περιόδου αναφοράς (ΠΡΙΝ την εγκατάσταση, Σεπ. 2024 -- Φεβ. 2025, Πάροχος: \textlatin{Protergia}) και της περιόδου βελτιστοποίησης (ΜΕΤΑ την εγκατάσταση, Μάι. 2025 -- Νοε. 2025, Πάροχος: \textlatin{NRG}). Ο Πίνακας Β.1 αποτυπώνει τη μη σταθμισμένη σύγκριση σε επίπεδο λογαριασμών, ενώ ο κύριος δείκτης 19,7\% της εργασίας προκύπτει από ξεχωριστή σταθμισμένη κανονικοποίηση μόνο για ημέρες πλήρους σχολικής λειτουργίας.

Η αλλαγή παρόχου μεταξύ των δύο περιόδων (\textlatin{Protergia} $\rightarrow$ \textlatin{NRG}) αποτελεί περιορισμό της σύγκρισης, καθώς ενδέχεται να επηρεάζει τις χρεώσεις δικτύου. Ωστόσο, η τιμή ενέργειας (\euro/\textlatin{kWh}) παρέμεινε σταθερή στα 0,162~\euro/\textlatin{kWh} και στις δύο περιόδους, διασφαλίζοντας συγκρισιμότητα στο επίπεδο κατανάλωσης.

\begin{table}[htbp]
\centering
\caption{Μη σταθμισμένη συγκριτική αποτύπωση βασικών ενεργειακών μετρικών σε επίπεδο λογαριασμών.}
\label{tab:appendix_energy_summary}
\begin{tabular}{@{}p{4.4cm}p{4.2cm}p{4.6cm}p{3.2cm}@{}}
\toprule
\textbf{Μετρική} & \textbf{Περίοδος Αναφοράς (ΠΡΙΝ)} & \textbf{Περίοδος Ψηφιακού Διδύμου (ΜΕΤΑ)} & \textbf{Ποσοστιαία Αλλαγή} \\
\midrule
\textbf{Περίοδος Μέτρησης} & Σεπ. 2024 -- Φεβ. 2025 & Μάι. 2025 -- Νοε. 2025 & -- \\
\textbf{Σύνολο Ημερών} & 122 τιμολογημένες ημέρες & 212 τιμολογημένες/έγκυρες ημέρες & -- \\
\textbf{Μέση Τιμή Ενέργειας} & 0,162 \euro/\textlatin{kWh} & 0,162 \euro/\textlatin{kWh} & 0,0\% \\
\textbf{Μέση Ημερήσια Κατανάλωση} & \textbf{201 \textlatin{kWh}/ημέρα} & \textbf{179 \textlatin{kWh}/ημέρα} & \textbf{-10,9\%*} \\
\bottomrule
\end{tabular}
\end{table}

\noindent Η επίδραση της αποκοπής των ``κρυφών φορτίων'' (\textlatin{phantom loads}) αποτυπώνεται σαφέστερα στη σύγκριση ίδιων μηνών (\textlatin{Year-over-Year}) του Πίνακα Β.2, με την ισχυρότερη βελτίωση να εμφανίζεται στον Νοέμβριο.

\begin{table}[htbp]
\centering
\caption{Αναλυτική μηνιαία σύγκριση κατανάλωσης (\textlatin{YoY} και \textlatin{vs Baseline}).}
\label{tab:appendix_energy_yoy}
\begin{tabular}{@{}lcccc@{}}
\toprule
\textbf{Μήνας} & \textbf{\textlatin{kWh}/ημέρα 2024} & \textbf{\textlatin{kWh}/ημέρα 2025} & \textlatin{YoY} \textbf{Μεταβολή} & \textbf{\textlatin{vs Baseline (223)}} \\
\midrule
Σεπτέμβριος & 163 & 197 & +20,9\% & -11,7\% \\
Οκτώβριος & 134 & 145 & +8,2\% & -35,0\% \\
Νοέμβριος & 252 & 170 & -32,5\% & -23,8\% \\
\bottomrule
\end{tabular}
\end{table}

\noindent Η στήλη ``\textlatin{YoY} Μεταβολή'' συγκρίνει τον ίδιο μήνα μεταξύ 2024 και 2025, αντανακλώντας εποχιακές και λειτουργικές συνθήκες. Η στήλη ``\textlatin{vs Baseline}'' συγκρίνει κάθε μήνα με τη σταθμισμένη μέση ημερήσια κατανάλωση της περιόδου αναφοράς (223 \textlatin{kWh}/ημέρα, Σεπ. 2024 -- Φεβ. 2025). Και οι δύο μετρικές είναι ορθές και απαντούν σε διαφορετικά ερευνητικά ερωτήματα (βλ. §6.2).


\noindent Η τιμή \texttt{-10,9\%} του Πίνακα Β.1 προκύπτει από μη σταθμισμένη σύγκριση συνολικών λογαριασμών. Η τιμή \texttt{-19,7\%} που αναφέρεται στο κυρίως κείμενο προκύπτει από σταθμισμένη κανονικοποίηση μόνο για ημέρες πλήρους σχολικής λειτουργίας --- βλ. \S~6.2 για αναλυτική επεξήγηση. Οι 212 ημέρες της στήλης ``ΜΕΤΑ'' αντιστοιχούν σε τιμολογημένες/έγκυρες ημέρες μέτρησης και όχι στο σύνολο των ημερολογιακών ημερών του διαστήματος Μαΐου--Νοεμβρίου.
